% =====================================================================
\BookChapter[Choosing a Language]{Choosing a Language and Writing Code Others Can Trust}
\label{ch:choosing}

\begin{researchquestion}
Which language should you use for a given task, and how do you make the result trustworthy?
\end{researchquestion}
% =====================================================================

You have now seen a dozen languages. This final chapter steps back and asks two practical
questions. Given an engineering task, which language should you use? And whatever the
language, how do you write code that you, your colleagues and your clients can trust?

\LearningOutcomes{Choose a language for a task by asking who will use the code, where it must run and how long it must live; learn a new language systematically; keep code under version control; test a calculation against known answers and limiting cases; document code so that results can be reproduced; and use AI coding assistants without giving up responsibility for the result.}

\section{Questions to ask before choosing}

There is no best language, only a best language \emph{for a given task, team and
situation}. The following questions usually decide the matter.

\begin{enumerate}
\item \textbf{What already exists?} If your company's simulation tool is scripted in \py,
  or your laboratory's instruments come with a LabVIEW driver, or your supervisor's
  solver is in Fortran, that is usually the strongest argument. Reusing tested code beats
  writing new code.
\item \textbf{Who will use and maintain it?} Code outlives its author's interest in it. A
  \py script can be read by most engineers; a clever C++ template library may be readable
  only by you. Write for the person who will inherit your code.
\item \textbf{Where must it run?} A web page needs JavaScript. A microcontroller needs C,
  C++ or Rust. A PLC needs IEC 61131-3 languages. A cluster job can use almost anything,
  but Fortran, C++ and \py dominate. A colleague without a MATLAB licence cannot run your
  MATLAB code.
\item \textbf{How fast must it be?} Most engineering scripts run in seconds, and speed
  does not matter. If it does, first find out \emph{where} the time is spent (profiling),
  then vectorise, use a library, or move the hot loop into a compiled language
  (Figure~\ref{fig:benchmark}).
\item \textbf{How safe must it be?} Software that controls a vehicle, a medical device or
  a chemical plant is subject to standards (ISO 26262, IEC 62304, IEC 61508) that restrict
  languages and require rigorous verification. There, C with MISRA rules, Ada/SPARK, Rust,
  certified code generators (Simulink, SCADE) and PLC languages are typical.
\item \textbf{What libraries do you need?} Statistics: R or \py. Machine learning: \py.
  Control design: MATLAB or \py. Differential equations at scale: Julia, \py, C++ or
  Fortran. Symbolic derivations: Wolfram, SymPy or Maple.
\item \textbf{How long will it live?} A throw-away analysis can be a notebook. A tool used
  by a hundred people for ten years needs a language with a stable future, tests,
  documentation and version control.
\end{enumerate}

\section{A rough decision guide}

Table~\ref{tab:ch16:guide} gives a first suggestion for some common tasks. Treat it as a
starting point for the questions above, not as a rule.

\begin{table}[htbp]
\centering\small
\caption[{A rough decision guide}]{A rough decision guide. The right-hand column names the
languages to consider first; local habits and existing code often decide otherwise.}
\label{tab:ch16:guide}
\begin{tabularx}{\linewidth}{@{}>{\raggedright\arraybackslash}p{5.6cm}>{\raggedright\arraybackslash}X@{}}
\toprule
\textbf{If you need to \ldots} & \textbf{consider first \ldots}\\
\midrule
analyse data, automate a task, prototype a model, drive other software & \py\\
do statistics, designed experiments, reliability analysis & R (or \py)\\
design control systems, process signals, model-based design & MATLAB/Simulink (or \py)\\
derive formulas, solve equations symbolically & Wolfram Language, SymPy, Maple\\
write a fast numerical solver from scratch & Julia, modern Fortran, C++\\
extend an existing large simulation code & the language it is written in\\
program a microcontroller or write firmware & C, C++, Rust\\
program an industrial controller (PLC) & Structured Text, Ladder Diagram\\
design digital hardware & VHDL, Verilog\\
build an interactive web tool or dashboard & JavaScript/TypeScript (or \py with Streamlit/Dash, R with Shiny)\\
build a Windows desktop application or a CAD plug-in & C\#\\
build an Android or iOS app & Kotlin, Swift\\
query or manage a database & SQL\\
automate a cluster workflow & Bash (plus \py)\\
model a multi-domain physical system & Modelica, Simscape, ModelingToolkit\\
\bottomrule
\end{tabularx}
\end{table}

\section{How to learn a new language quickly}

Sooner or later you will need a language you do not know. A systematic approach:
\begin{enumerate}
\item \textbf{Get ``Hello, world'' running.} Install the tools and learn the
  edit--run cycle. This is often the hardest step.
\item \textbf{Map the core features} using the checklist from Chapter~\ref{ch:intro}:
  types, variables, operators, \code{if}, loops, functions, arrays, I/O, modules, errors.
  Table~\ref{tab:syntax} and the reference guides at the back help.
\item \textbf{Translate a program you already understand}, such as the drag or Rosetta
  programs of these notes. Because you know what the output must be, every difference
  teaches you something.
\item \textbf{Read idiomatic code} --- the standard library, well-known packages, the
  official tutorial --- to learn how experienced users write the language, not just what
  compiles.
\item \textbf{Learn the tools}: the package manager, the debugger, the test framework, the
  formatter.
\end{enumerate}

\section{Version control with Git}

\term[version control]{Version control} records every change to a set of files, who made it and why, and
lets you go back to any earlier state. It is essential for code, and useful for any text
files (\LaTeX\ reports, data-processing scripts, configuration). \textbf{Git} is the
standard tool; \textbf{GitHub}, \textbf{GitLab} and similar services host Git
repositories online and support collaboration.
\begin{lstlisting}[style=shell]
$ git init                         # start tracking the current folder
$ git add drag.py                  # stage a file for the next snapshot
$ git commit -m "Euler solver for drag problem"   # record the snapshot
$ git log --oneline                # list the history
$ git diff                         # what changed since the last commit?
$ git checkout -b smaller-dt       # create a branch to try an idea
$ git push origin main             # upload to GitHub/GitLab
\end{lstlisting}
Commit often, with messages that explain \emph{why} something changed. Never commit
passwords or large generated data files. Graphical tools such as GitHub Desktop and the Git
support built into VS Code make all of this easier.

\section{Testing}

A calculation you have not checked is a guess. Chapter~\ref{ch:python} introduced
automated tests; here are the kinds of checks engineers should build into their programs,
in any language:
\begin{itemize}
\item \textbf{Known solutions}: compare with an exact solution (as we did with
  $\tanh$), a textbook example, or a hand calculation.
\item \textbf{Limiting cases}: with zero drag the speed must be $gt$; with zero load the
  deflection must be zero.
\item \textbf{Conservation laws and symmetries}: energy in the MD simulation, mass and
  energy balances in process models, symmetric stiffness matrices.
\item \textbf{Convergence}: halving the time step or refining the mesh should change the
  result by the amount the method's order predicts.
\item \textbf{Units and dimensions}: check every formula; use unit-aware libraries where
  possible.
\item \textbf{Independent implementation}: as in Chapter~\ref{ch:rosetta}, a second
  implementation in another language or by another person catches errors that a single
  author does not see.
\item \textbf{Regression tests}: once a result is verified, store it and check that future
  versions of the code still reproduce it.
\end{itemize}
Every language has test frameworks: \code{pytest} (\py), \code{testthat} (R), MATLAB's
\code{matlab.unittest}, JUnit (Java), xUnit/NUnit (C\#), GoogleTest and Catch2 (C++),
\code{pFUnit} and \code{test-drive} (Fortran), Julia's \code{Test}, Rust's built-in
\code{cargo test}, and Jest or Vitest (JavaScript).

\section{Documentation and reproducibility}

Good engineering code is \emph{documented}: every function says what it does, what its
inputs and outputs are, and \emph{in which units}; a README file explains how to install
and run the program; and assumptions and limitations are written down. Good engineering
results are \emph{reproducible}: someone else, or you next year, can regenerate every
number and figure from the raw data and the code. That requires recording the versions of
the language and of every package (\file{requirements.txt}, \file{environment.yml},
\file{renv.lock}, \file{Project.toml}, \file{Cargo.lock}), fixing random seeds, and keeping
code, data and results under version control. Publishing code and data alongside a report
or paper, in a public repository with a clear licence, is increasingly expected.

\section{AI coding assistants}
\label{sec:ai-assistants}

Since about 2022, AI assistants based on large language models have become part of
everyday programming. They can explain unfamiliar code, translate a program from MATLAB to
\py, write a first version of a function from a description, suggest tests, and find bugs.
They know every language in these notes. Used well, they make you much more productive and
are an excellent tutor for learning a new language.

They can also be confidently wrong. An assistant may invent a function that does not
exist, use a formula with a wrong constant, mix up 0- and 1-based indices when translating,
or produce code that runs and gives a plausible but incorrect answer --- the most dangerous
kind of error (Chapter~\ref{ch:howcoderuns}). The engineering responsibility for a result
remains with the engineer who signs it. Some rules for using AI assistants:
\begin{enumerate}
\item \textbf{Understand every line} you use. If you could not explain it to a colleague,
  you cannot vouch for it. This is why reading code in many languages --- the purpose of
  these notes --- matters more, not less, in the age of AI.
\item \textbf{Test generated code} with the same checks as your own: known solutions,
  limiting cases, units.
\item \textbf{Give precise specifications}: inputs, outputs, units, the method to use, and
  an example with the expected answer.
\item \textbf{Follow your course and employer rules} on what may be shared with external
  services (confidential data, proprietary code) and on how AI use must be acknowledged in
  assessed work.
\end{enumerate}

\section{Software licences}

When you use or share code, its \term{licence} decides what you may do with it. Open-source
licences fall into two broad groups. \emph{Permissive} licences (MIT, BSD, Apache 2.0) allow
almost any use, including in commercial products, provided the copyright notice is kept.
\emph{Copyleft} licences (GPL, LGPL) require that modified versions you distribute are
released under the same licence. Data and documents are often shared under Creative Commons
licences (CC BY 4.0, for example, requires attribution). Commercial software (MATLAB,
Mathematica, many simulation packages) comes with licence terms that may restrict who can
run it and for what purpose. Check before you build on someone else's code.

\section{A last word}

Choosing a language is an engineering decision, and it deserves the same care as choosing
a material or a safety factor. Learn \py well, learn to read the others, and check what
you compute. The Epilogue that follows looks at where these languages came from and
where they may be going.

\begin{exercises}
\item For each of the following tasks, choose a language and justify your choice in two or
  three sentences: (a) logging and plotting data from a temperature sensor on a Raspberry
  Pi; (b) a cost estimate that the finance department will check; (c) a controller for a
  conveyor belt; (d) a 3D simulation of airflow around a building on a cluster; (e) an
  interactive web page that lets clients size a heat exchanger; (f) a statistical analysis
  of fatigue-test results.
\item Create a GitHub repository for your ENG 101 code. Add a README that explains how to
  run each program, a \file{requirements.txt}, and a test for at least one function.
\item Ask an AI assistant to translate \file{drag.m} into \py. Check the translation line
  by line. Did it handle the indices and the time vector correctly? Did it produce the
  same table?
\item Write a test that checks the convergence order of the Euler method for the drag
  problem: the error at $t = 6\unit{s}$ should halve when $\Delta t$ is halved.
\item Read the licence of NumPy and of one commercial program you use. What may you do
  with each, and what may you not?
\end{exercises}
